\documentclass[11pt]{article}

\usepackage[T1]{fontenc}
\usepackage[utf8]{inputenc}
\usepackage[spanish,es-nodecimaldot]{babel}
\usepackage{amsmath,amssymb,mathtools}
\usepackage{enumitem}
\usepackage[a4paper,margin=2.5cm]{geometry}
\usepackage[hidelinks]{hyperref}

\newcommand{\Cstar}{C^*}
\newcommand{\GL}{\mathrm{GL}}
\newcommand{\Cb}{C_b}
\newcommand{\Czero}{C_0}
\newcommand{\one}{\mathbf{1}}

\newenvironment{definicion}[1]
  {\par\medskip\noindent\textbf{Definición #1.}\ }
  {\par\medskip}

\title{Definiciones de la Parte 1 de $C^*$-álgebras: Teoría espectral en álgebras de Banach}
\author{}
\date{}

\begin{document}
\maketitle

\section*{Índice de definiciones}
\begin{itemize}[leftmargin=*]
  \item \textbf{Definición 1.1.} Álgebra de Banach; unidad y conmutatividad. \hfill (p. 1)
  \item \textbf{Definición 1.4.} Involución y álgebra de Banach involutiva. \hfill (p. 1)
  \item \textbf{Definición 1.7.} Grupo de elementos invertibles $\GL(A)$. \hfill (p. 2)
  \item \textbf{Definición 1.11.} Espectro de un elemento. \hfill (p. 3)
  \item \textbf{Definición 1.13.} Conjunto resolvente y función resolvente. \hfill (p. 3)
  \item \textbf{Definición 1.17.} Carácter. \hfill (p. 4)
  \item \textbf{Definición 1.19.} Espectro de Gelfand o espacio de caracteres. \hfill (p. 4)
  \item \textbf{Definición 1.25.} Radio espectral. \hfill (p. 6)
  \item \textbf{Definición 1.27.} Transformada de Gelfand. \hfill (p. 9)
  \item \textbf{Definición 1.29.} Ideal modular. \hfill (p. 9)
\end{itemize}

\section*{Definiciones}

% Fuente: Parte1.pdf, p. 1
\begin{definicion}{1.1}
Un \emph{álgebra de Banach} es un álgebra asociativa $A$ sobre $\mathbb{C}$ con una norma tal que
\begin{itemize}[leftmargin=2em]
  \item $(A,+,\lVert\cdot\rVert)$ es un espacio de Banach;
  \item para todo $x,y\in A$,
  \[
    \lVert xy\rVert\leq \lVert x\rVert\,\lVert y\rVert.
  \]
\end{itemize}
Además decimos que $A$ es:
\begin{itemize}[leftmargin=2em]
  \item \emph{con unidad} si tiene un neutro multiplicativo $1_A$ que satisface
  \[
    \lVert 1_A\rVert=1;
  \]
  \item \emph{conmutativa} si
  \[
    xy=yx
    \qquad\text{para todo }x,y\in A.
  \]
\end{itemize}
\end{definicion}

% Fuente: Parte1.pdf, p. 1
\begin{definicion}{1.4}
Sea $A$ un álgebra compleja. Una \emph{involución} en $A$ es un mapeo
\[
  *:A\longrightarrow A
\]
tal que
\[
  (\lambda a+b)^*=\overline{\lambda}\,a^*+b^*,
  \qquad
  (a^*)^*=a,
  \qquad
  (ab)^*=b^*a^*.
\]
Si $A$ está equipada con una involución, diremos que $A$ es un \emph{álgebra involutiva}.

Si $A$ es además un álgebra de Banach tal que
\[
  \lVert a^*\rVert=\lVert a\rVert,
\]
diremos que $A$ es un \emph{álgebra de Banach involutiva}.
\end{definicion}

% Fuente: Parte1.pdf, p. 2
\begin{definicion}{1.7}
Sea $A$ un álgebra de Banach con unidad. Definimos
\[
  \GL(A)=\{\text{elementos invertibles de }A\}.
\]
\end{definicion}

% Fuente: Parte1.pdf, p. 3
\begin{definicion}{1.11}
Sea $A$ un álgebra de Banach con unidad y $x\in A$. El \emph{espectro} de $x$ se define como
\[
  \sigma(x)=\sigma_A(x)
  =\{\lambda\in\mathbb{C}:\lambda 1_A-x\notin\GL(A)\}.
\]
\end{definicion}

% Fuente: Parte1.pdf, p. 3
\begin{definicion}{1.13}
Sea $A$ un álgebra de Banach y $x\in A$. Definimos el \emph{conjunto resolvente} de $x$ como
\[
  \rho(x)=\mathbb{C}\setminus\sigma(x),
\]
y la \emph{función resolvente}
\[
  R_x:\rho(x)\longrightarrow A,
  \qquad
  R_x(z)=(z1_A-x)^{-1}.
\]
\end{definicion}

% Fuente: Parte1.pdf, p. 4
\begin{definicion}{1.17}
Sea $A$ un álgebra de Banach. Un \emph{carácter} en $A$ es un funcional lineal multiplicativo no nulo
\[
  \varphi:A\longrightarrow\mathbb{C}.
\]
\end{definicion}

% Fuente: Parte1.pdf, p. 4
\begin{definicion}{1.19}
Sea $A$ un álgebra de Banach conmutativa con unidad. Definimos el \emph{espectro} (o \emph{espacio de caracteres}, o \emph{espectro de Gelfand}) como
\[
  \Omega(A)=\{\text{caracteres }A\to\mathbb{C}\}\subseteq A^*.
\]
Equipamos $\Omega(A)$ con la topología $w^*$.
\end{definicion}

% Fuente: Parte1.pdf, p. 6
\begin{definicion}{1.25}
Sea $A$ un álgebra de Banach con unidad y $x\in A$. Definimos el \emph{radio espectral} como
\[
  r(x)=\max\{\lvert\lambda\rvert:\lambda\in\sigma(x)\}.
\]
\end{definicion}

% Fuente: Parte1.pdf, p. 9
\begin{definicion}{1.27}
Sea $A$ un álgebra de Banach conmutativa con unidad. Se define la \emph{transformada de Gelfand} como el mapeo
\[
  \Lambda:A\longrightarrow C(\Omega(A)),
  \qquad
  x\longmapsto \widehat{x},
  \qquad
  \widehat{x}(\varphi):=\varphi(x).
\]
\end{definicion}

% Fuente: Parte1.pdf, p. 9
\begin{definicion}{1.29}
Un ideal $I$ de $A$ es \emph{modular} si existe $u\in A$ tal que
\[
  a-ua\in I,
  \qquad
  a-au\in I,
  \qquad
  \forall a\in A.
\]
\end{definicion}

\end{document}
